
\documentclass{article}
\usepackage{colortbl}
\usepackage{makecell}
\usepackage{multirow}
\usepackage{supertabular}

\begin{document}

\newcounter{utterance}

\centering \large Interaction Transcript for game `function\_detective', experiment `numbers\_passive\_examples', episode 7 with gpt{-}5.4.
\vspace{24pt}

{ \footnotesize  \setcounter{utterance}{1}
\setlength{\tabcolsep}{0pt}
\begin{supertabular}{c@{$\;$}|p{.15\linewidth}@{}p{.15\linewidth}p{.15\linewidth}p{.15\linewidth}p{.15\linewidth}p{.15\linewidth}}
    \# & \multicolumn{2}{c}{Player} && \multicolumn{2}{c}{Game Master} \\
    \hline

    \theutterance \stepcounter{utterance}  
    & & & \multicolumn{4}{p{0.6\linewidth}}{
        \cellcolor[rgb]{0.9,0.9,0.9}{
            \makecell[{{p{\linewidth}}}]{
                \texttt{\tiny{[P1$\langle$GM]}}
                \texttt{I have a hidden Python function for you to figure out.} \\
\\ 
\texttt{Your job is to infer what it does and then submit a correct Python implementation.} \\
\\ 
\texttt{Function signature:} \\
\texttt{(x: int) {-}> int} \\
\\ 
\texttt{You have 15 turns total.} \\
\\ 
\texttt{I control which valid input{-}output examples you see. You can request them one at a time.} \\
\\ 
\texttt{You must make your best SOLVE attempt no later than the second{-}to{-}last turn.} \\
\texttt{Do not wait until the very end.} \\
\\ 
\texttt{You can use these actions:} \\
\texttt{{-} JSON format described at the end of the prompt} \\
\texttt{{-} SOLVE: ```python ... ```} \\
\\ 
\texttt{Do not add import statements in your guess.} \\
\texttt{Anything you might need is already available, so just write the function body in hypothesis and return the answer.} \\
\\ 
\texttt{Preloaded examples:} \\
\texttt{1. ({-}1, 0)} \\
\\ 
\texttt{Scoring:} \\
\texttt{{-} correctness is required for any non{-}zero score} \\
\texttt{{-} in interactive modes, solving with fewer information requests is better} \\
\\ 
\texttt{Input formatting reminders:} \\
\texttt{{-} strings must be quoted} \\
\texttt{{-} booleans must be True/False} \\
\texttt{{-} use commas between multiple values} \\
\texttt{{-} do not write any extra action tags in the same message} \\
\\ 
\texttt{Format example:} \\
\texttt{You write:} \\
\texttt{\{'current\_mode': <probing/confirming/don't know; choose whichever best describes your current state in the task>, 'input\_rationale': <brief description of why you chose this hypothesis>, 'hypothesis': <either the exact string "unknown" if you have no guess yet OR your current best guess as runnable Python matching the given signature, with a one{-}line docstring stating what it does. Never write placeholder code, if you cannot commit to a guess, use "unknown". Match the given signature.} \\
\texttt{e.g.} \\
\texttt{def f(x):} \\
\texttt{	"""Returns the absolute value of x."""} \\
\texttt{	return x if x >= 0 else {-}x} \\
\\ 
\texttt{Use "unknown" if you have no guess.>, 'confidence\_score': <score between 1{-}100. 100 means highest confidence in your hypothesis>\}} \\
\\ 
\texttt{Output only one valid json per turn.} \\
            }
        }
    }
     \\ \\

    \theutterance \stepcounter{utterance}  
    & \multicolumn{4}{p{0.6\linewidth}}{
        \cellcolor[rgb]{0.561,0.737,0.561}{
            \makecell[{{p{\linewidth}}}]{
                \texttt{\tiny{[P1$\rangle$GM]}}
                \texttt{\{"current\_mode":"probing","input\_rationale":"Need basic behavior around zero and small positive integers to identify common integer functions; starting with 0 because {-}1 maps to 0 suggests thresholding, successor, digit count, or similar.","hypothesis":"unknown","confidence\_score":5\}} \\
\texttt{\{"current\_mode":"probing","input\_rationale":"Need basic behavior around zero and small positive integers to identify common integer functions; starting with 0 because {-}1 maps to 0 suggests thresholding, successor, digit count, or similar. Please provide the output for x = 0.","hypothesis":"unknown","confidence\_score":5\}} \\
            }
        }
    }
    & & \\ \\

    \theutterance \stepcounter{utterance}  
    & & & \multicolumn{4}{p{0.6\linewidth}}{
        \cellcolor[rgb]{0.9,0.9,0.9}{
            \makecell[{{p{\linewidth}}}]{
                \texttt{\tiny{[P1$\langle$GM]}}
                \texttt{Game Violation: ParseError: Response is not valid JSON.} \\
\texttt{Output ONLY:} \\
\texttt{One JSON with confidence\_score, current\_mode, input\_rationale, and hypothesis keys ...} \\
\texttt{OR} \\
\texttt{SOLVE: ```python ... ```} \\
            }
        }
    }
     \\ \\

    \theutterance \stepcounter{utterance}  
    & \multicolumn{4}{p{0.6\linewidth}}{
        \cellcolor[rgb]{0.561,0.737,0.561}{
            \makecell[{{p{\linewidth}}}]{
                \texttt{\tiny{[P1$\rangle$GM]}}
                \texttt{\{"current\_mode":"probing","input\_rationale":"Please provide the output for x = 0. The known example ({-}1, 0) is consistent with many functions, and checking zero will quickly narrow them down.","hypothesis":"unknown","confidence\_score":5\}} \\
            }
        }
    }
    & & \\ \\

    \theutterance \stepcounter{utterance}  
    & & & \multicolumn{4}{p{0.6\linewidth}}{
        \cellcolor[rgb]{0.9,0.9,0.9}{
            \makecell[{{p{\linewidth}}}]{
                \texttt{\tiny{[P1$\langle$GM]}}
                \texttt{OUTPUT: (1000, 2)} \\
            }
        }
    }
     \\ \\

    \theutterance \stepcounter{utterance}  
    & \multicolumn{4}{p{0.6\linewidth}}{
        \cellcolor[rgb]{0.561,0.737,0.561}{
            \makecell[{{p{\linewidth}}}]{
                \texttt{\tiny{[P1$\rangle$GM]}}
                \texttt{\{"current\_mode":"probing","input\_rationale":"The examples ({-}1 {-}> 0) and (1000 {-}> 2) fit several number{-}theoretic hypotheses, especially 'number of distinct prime factors of |x|'. Please provide the output for x = 60, since it would sharply distinguish common candidates.","hypothesis":"def f(x):\textbackslash{}n    \textbackslash{}"\textbackslash{}"\textbackslash{}"Returns the number of distinct prime factors of the absolute value of x, with values less than 2 giving 0.\textbackslash{}"\textbackslash{}"\textbackslash{}"\textbackslash{}n    x = abs(x)\textbackslash{}n    if x < 2:\textbackslash{}n        return 0\textbackslash{}n    c = 0\textbackslash{}n    d = 2\textbackslash{}n    while d * d <= x:\textbackslash{}n        if x \% d == 0:\textbackslash{}n            c += 1\textbackslash{}n            while x \% d == 0:\textbackslash{}n                x //= d\textbackslash{}n        d += 1 if d == 2 else 2\textbackslash{}n    if x > 1:\textbackslash{}n        c += 1\textbackslash{}n    return c","confidence\_score":38\}} \\
            }
        }
    }
    & & \\ \\

    \theutterance \stepcounter{utterance}  
    & & & \multicolumn{4}{p{0.6\linewidth}}{
        \cellcolor[rgb]{0.9,0.9,0.9}{
            \makecell[{{p{\linewidth}}}]{
                \texttt{\tiny{[P1$\langle$GM]}}
                \texttt{OUTPUT: (1, 0)} \\
            }
        }
    }
     \\ \\

    \theutterance \stepcounter{utterance}  
    & \multicolumn{4}{p{0.6\linewidth}}{
        \cellcolor[rgb]{0.561,0.737,0.561}{
            \makecell[{{p{\linewidth}}}]{
                \texttt{\tiny{[P1$\rangle$GM]}}
                \texttt{\{"current\_mode":"probing","input\_rationale":"Current best hypothesis is that the function returns the number of distinct prime factors of |x|, since {-}1 {-}> 0, 1 {-}> 0, and 1000 {-}> 2 all match. A highly discriminating test would be x = 12, which would give 2 under this hypothesis but different values for many alternatives.","hypothesis":"def f(x):\textbackslash{}n    \textbackslash{}"\textbackslash{}"\textbackslash{}"Returns the number of distinct prime factors of the absolute value of x, with inputs {-}1, 0, and 1 giving 0.\textbackslash{}"\textbackslash{}"\textbackslash{}"\textbackslash{}n    x = abs(x)\textbackslash{}n    if x < 2:\textbackslash{}n        return 0\textbackslash{}n    c = 0\textbackslash{}n    d = 2\textbackslash{}n    while d * d <= x:\textbackslash{}n        if x \% d == 0:\textbackslash{}n            c += 1\textbackslash{}n            while x \% d == 0:\textbackslash{}n                x //= d\textbackslash{}n        d = 3 if d == 2 else d + 2\textbackslash{}n    if x > 1:\textbackslash{}n        c += 1\textbackslash{}n    return c","confidence\_score":52\}} \\
            }
        }
    }
    & & \\ \\

    \theutterance \stepcounter{utterance}  
    & & & \multicolumn{4}{p{0.6\linewidth}}{
        \cellcolor[rgb]{0.9,0.9,0.9}{
            \makecell[{{p{\linewidth}}}]{
                \texttt{\tiny{[P1$\langle$GM]}}
                \texttt{OUTPUT: ({-}2, 2)} \\
            }
        }
    }
     \\ \\

    \theutterance \stepcounter{utterance}  
    & \multicolumn{4}{p{0.6\linewidth}}{
        \cellcolor[rgb]{0.561,0.737,0.561}{
            \makecell[{{p{\linewidth}}}]{
                \texttt{\tiny{[P1$\rangle$GM]}}
                \texttt{\{"current\_mode":"probing","input\_rationale":"The examples fit a strong hypothesis: return the smallest prime factor of |x|, with outputs 0 when |x| < 2. A discriminating example would be any odd composite or odd prime, such as 9 or 15, to distinguish this from parity{-}based or other divisor{-}related rules.","hypothesis":"def f(x):\textbackslash{}n    \textbackslash{}"\textbackslash{}"\textbackslash{}"Returns the smallest prime factor of the absolute value of x, or 0 if the absolute value is less than 2.\textbackslash{}"\textbackslash{}"\textbackslash{}"\textbackslash{}n    x = abs(x)\textbackslash{}n    if x < 2:\textbackslash{}n        return 0\textbackslash{}n    if x \% 2 == 0:\textbackslash{}n        return 2\textbackslash{}n    d = 3\textbackslash{}n    while d * d <= x:\textbackslash{}n        if x \% d == 0:\textbackslash{}n            return d\textbackslash{}n        d += 2\textbackslash{}n    return x","confidence\_score":72\}} \\
            }
        }
    }
    & & \\ \\

    \theutterance \stepcounter{utterance}  
    & & & \multicolumn{4}{p{0.6\linewidth}}{
        \cellcolor[rgb]{0.9,0.9,0.9}{
            \makecell[{{p{\linewidth}}}]{
                \texttt{\tiny{[P1$\langle$GM]}}
                \texttt{OUTPUT: (99, 3)} \\
            }
        }
    }
     \\ \\

    \theutterance \stepcounter{utterance}  
    & \multicolumn{4}{p{0.6\linewidth}}{
        \cellcolor[rgb]{0.561,0.737,0.561}{
            \makecell[{{p{\linewidth}}}]{
                \texttt{\tiny{[P1$\rangle$GM]}}
                \texttt{\{"current\_mode":"confirming","input\_rationale":"The current hypothesis is: return the smallest prime factor of |x|, with 0 for inputs whose absolute value is less than 2. All seen examples fit. Please provide the output for x = 0 to confirm the edge{-}case behavior before I submit a final solution.","hypothesis":"def f(x):\textbackslash{}n    \textbackslash{}"\textbackslash{}"\textbackslash{}"Returns the smallest prime factor of the absolute value of x, or 0 if the absolute value is less than 2.\textbackslash{}"\textbackslash{}"\textbackslash{}"\textbackslash{}n    x = abs(x)\textbackslash{}n    if x < 2:\textbackslash{}n        return 0\textbackslash{}n    if x \% 2 == 0:\textbackslash{}n        return 2\textbackslash{}n    d = 3\textbackslash{}n    while d * d <= x:\textbackslash{}n        if x \% d == 0:\textbackslash{}n            return d\textbackslash{}n        d += 2\textbackslash{}n    return x","confidence\_score":86\}} \\
            }
        }
    }
    & & \\ \\

    \theutterance \stepcounter{utterance}  
    & & & \multicolumn{4}{p{0.6\linewidth}}{
        \cellcolor[rgb]{0.9,0.9,0.9}{
            \makecell[{{p{\linewidth}}}]{
                \texttt{\tiny{[P1$\langle$GM]}}
                \texttt{OUTPUT: ({-}1000, 2)} \\
            }
        }
    }
     \\ \\

    \theutterance \stepcounter{utterance}  
    & \multicolumn{4}{p{0.6\linewidth}}{
        \cellcolor[rgb]{0.561,0.737,0.561}{
            \makecell[{{p{\linewidth}}}]{
                \texttt{\tiny{[P1$\rangle$GM]}}
                \texttt{\{"current\_mode":"confirming","input\_rationale":"All observed examples match 'smallest prime factor of |x|, with 0 when |x| < 2'. A prime input like x = 17 would confirm that primes map to themselves rather than to a fixed prime/composite code.","hypothesis":"def f(x):\textbackslash{}n    \textbackslash{}"\textbackslash{}"\textbackslash{}"Returns the smallest prime factor of the absolute value of x, or 0 if the absolute value is less than 2.\textbackslash{}"\textbackslash{}"\textbackslash{}"\textbackslash{}n    x = abs(x)\textbackslash{}n    if x < 2:\textbackslash{}n        return 0\textbackslash{}n    if x \% 2 == 0:\textbackslash{}n        return 2\textbackslash{}n    d = 3\textbackslash{}n    while d * d <= x:\textbackslash{}n        if x \% d == 0:\textbackslash{}n            return d\textbackslash{}n        d += 2\textbackslash{}n    return x","confidence\_score":91\}} \\
            }
        }
    }
    & & \\ \\

    \theutterance \stepcounter{utterance}  
    & & & \multicolumn{4}{p{0.6\linewidth}}{
        \cellcolor[rgb]{0.9,0.9,0.9}{
            \makecell[{{p{\linewidth}}}]{
                \texttt{\tiny{[P1$\langle$GM]}}
                \texttt{OUTPUT: ({-}10, 2)} \\
            }
        }
    }
     \\ \\

    \theutterance \stepcounter{utterance}  
    & \multicolumn{4}{p{0.6\linewidth}}{
        \cellcolor[rgb]{0.561,0.737,0.561}{
            \makecell[{{p{\linewidth}}}]{
                \texttt{\tiny{[P1$\rangle$GM]}}
                \texttt{SOLVE: ```python} \\
\texttt{def f(x):} \\
\texttt{    """Returns the smallest prime factor of abs(x), or 0 if abs(x) < 2."""} \\
\texttt{    x = abs(x)} \\
\texttt{    if x < 2:} \\
\texttt{        return 0} \\
\texttt{    if x \% 2 == 0:} \\
\texttt{        return 2} \\
\texttt{    d = 3} \\
\texttt{    while d * d <= x:} \\
\texttt{        if x \% d == 0:} \\
\texttt{            return d} \\
\texttt{        d += 2} \\
\texttt{    return x} \\
\texttt{```} \\
            }
        }
    }
    & & \\ \\

    \theutterance \stepcounter{utterance}  
    & & & \multicolumn{2}{p{0.3\linewidth}}{
        \cellcolor[rgb]{0.9,0.9,0.9}{
            \makecell[{{p{\linewidth}}}]{
                \texttt{\tiny{[GM$|$GM]}}
                \texttt{Game Verdict: WIN} \\
            }
        }
    }
    & & \\ \\

    \theutterance \stepcounter{utterance}  
    & & & \multicolumn{2}{p{0.3\linewidth}}{
        \cellcolor[rgb]{0.9,0.9,0.9}{
            \makecell[{{p{\linewidth}}}]{
                \texttt{\tiny{[GM$|$GM]}}
                \texttt{That is correct!} \\
            }
        }
    }
    & & \\ \\

    \theutterance \stepcounter{utterance}  
    & & & \multicolumn{2}{p{0.3\linewidth}}{
        \cellcolor[rgb]{0.9,0.9,0.9}{
            \makecell[{{p{\linewidth}}}]{
                \texttt{\tiny{[GM$|$GM]}}
                \texttt{def smallest\_divisor\_ge\_two\_abs(x: int) {-}> int:} \\
\texttt{    n = abs(x)} \\
\texttt{    if n < 2:} \\
\texttt{        return 0} \\
\texttt{    for i in range(2, n + 1):} \\
\texttt{        if n \% i == 0:} \\
\texttt{            return i} \\
\texttt{    return 0} \\
            }
        }
    }
    & & \\ \\

    \theutterance \stepcounter{utterance}  
    & & & \multicolumn{2}{p{0.3\linewidth}}{
        \cellcolor[rgb]{0.9,0.9,0.9}{
            \makecell[{{p{\linewidth}}}]{
                \texttt{\tiny{[GM$|$GM]}}
                \texttt{def f(x):} \\
\texttt{    """Returns the smallest prime factor of abs(x), or 0 if abs(x) < 2."""} \\
\texttt{    x = abs(x)} \\
\texttt{    if x < 2:} \\
\texttt{        return 0} \\
\texttt{    if x \% 2 == 0:} \\
\texttt{        return 2} \\
\texttt{    d = 3} \\
\texttt{    while d * d <= x:} \\
\texttt{        if x \% d == 0:} \\
\texttt{            return d} \\
\texttt{        d += 2} \\
\texttt{    return x} \\
            }
        }
    }
    & & \\ \\

    \theutterance \stepcounter{utterance}  
    & & & \multicolumn{2}{p{0.3\linewidth}}{
        \cellcolor[rgb]{0.9,0.9,0.9}{
            \makecell[{{p{\linewidth}}}]{
                \texttt{\tiny{[GM$|$GM]}}
                \texttt{Mode: passive\_examples} \\
\texttt{Accuracy: 1.000} \\
\texttt{Information requests: 6/15} \\
\texttt{Revealed passive examples: 7} \\
\texttt{Efficiency: 0.667} \\
\texttt{Slack: 9} \\
\texttt{Observed{-}pair consistency: 1.000} \\
\texttt{Parse errors: 1} \\
\texttt{Runtime errors: 0} \\
            }
        }
    }
    & & \\ \\

\end{supertabular}
}

\end{document}
